% Modelo Beamer UFC original de Maurício Moreira Neto.
\ifdefined\Resumo\PassOptionsToClass{handout}{beamer}\fi
\documentclass{libs/ufc_format}
\usepackage{amsmath,amssymb,booktabs,tikz,circuitikz}
\usetikzlibrary{arrows.meta,positioning,calc,shapes.geometric}
% A geometria 4:3 e todos os elementos do tema vêm da classe ORIGINAL.
% Não substituir fundo, logo, cabeçalho, rodapé, fontes ou blocos do template.
\hypersetup{pdflang=pt-BR}
\usepackage{adjustbox}
\newcommand{\VDD}{V_{\mathrm{DD}}}
\newcommand{\GND}{\mathrm{GND}}
\newcommand{\fonte}[1]{\par\vspace{1mm}{\tiny\color{blue_theme}#1\par}}
\newcommand{\ideia}[1]{\begin{block}{Ideia-chave}#1\end{block}}
\newcommand{\pergunta}[1]{\begin{exampleblock}{Para discutir}#1\end{exampleblock}}
\newcommand{\duas}[2]{\begin{columns}[T,onlytextwidth]\begin{column}{.48\textwidth}#1\end{column}\begin{column}{.48\textwidth}#2\end{column}\end{columns}}
\newcommand{\topico}[2]{\begin{block}{#1}#2\end{block}}
\newcommand{\capa}[2]{\subtitle{#1\\[1mm]{\small #2}}\begin{frame}{}\maketitle\end{frame}}
% Redimensiona somente diagramas e tabelas que excedam a largura de texto.
\tikzset{caixa/.style={draw=blue_theme,fill=blue_theme!5,rounded corners=2pt,align=center,minimum height=9mm,font=\small},seta/.style={-{Latex},thick,blue_theme}}
\ctikzset{transistors/scale=.78}
\newcommand{\inversor}{%
\begin{circuitikz}[scale=.85,transform shape]
\node[pmos] (p) at (0,1.3) {};\node[nmos] (n) at (0,-1.1) {};
\draw (p.S) -- (0,2.2) node[above]{$\VDD$};
\draw (p.D) -- (n.D);\draw (0,.1) to[short,*-] (1.35,.1) node[right]{$Y$};
\draw (n.S) -- (0,-2) node[ground]{};
\draw (p.G) -- (-1.4,1.3) -- (-1.4,-1.1) -- (n.G);
\draw (-1.4,.1) to[short,*-] (-2.05,.1) node[left]{$A$};
\node[right] at (.25,1.3) {PMOS};\node[right] at (.25,-1.1) {NMOS};
\end{circuitikz}}
\newcommand{\nandcircuito}{%
\begin{circuitikz}[scale=.68,transform shape]
\node[pmos] (pa) at (-1.2,1.7) {};\node[pmos] (pb) at (1.4,1.7) {};
\node[nmos] (na) at (0,-.7) {};\node[nmos] (nb) at (0,-2.3) {};
\draw (pa.S) -- (-1.2,2.6) -- (1.4,2.6) -- (pb.S);\node[above] at (.1,2.6) {$\VDD$};
\draw (pa.D) -- (-1.2,.5) -- (1.4,.5) -- (pb.D);
\draw (0,.5) -- (na.D);\draw (na.S) -- (nb.D);\draw (nb.S) -- (0,-3.15) node[ground]{};
\draw (1.4,.5) to[short,*-] (2.15,.5) node[right]{$Y$};
\draw (pa.G) node[left]{$A$};\draw (pb.G) node[left]{$B$};\draw (na.G) node[left]{$A$};\draw (nb.G) node[left]{$B$};
\end{circuitikz}}
\newcommand{\norcircuito}{%
\begin{circuitikz}[scale=.68,transform shape]
\node[pmos] (pa) at (0,2) {};\node[pmos] (pb) at (0,.4) {};
\node[nmos] (na) at (-1.2,-1.8) {};\node[nmos] (nb) at (1.4,-1.8) {};
\draw (pa.S) -- (0,2.85) node[above]{$\VDD$};\draw (pa.D) -- (pb.S);\draw (pb.D) -- (0,-.6);
\draw (na.D) -- (-1.2,-.6) -- (1.4,-.6) -- (nb.D);
\draw (1.4,-.6) to[short,*-] (2.15,-.6) node[right]{$Y$};
\draw (na.S) -- (-1.2,-2.8) -- (1.4,-2.8) -- (nb.S);\draw (0,-2.8) node[ground]{};
\draw (pa.G) node[left]{$A$};\draw (pb.G) node[left]{$B$};\draw (na.G) node[left]{$A$};\draw (nb.G) node[left]{$B$};
\end{circuitikz}}
\newcommand{\aoi}{%
\begin{circuitikz}[scale=.64,transform shape]
\node[pmos] (pa) at (-1.3,2) {};\node[pmos] (pb) at (1.3,2) {};\node[pmos] (pc) at (0,.4) {};
\draw (pa.S) -- (-1.3,2.8) -- (1.3,2.8) -- (pb.S);\node[above] at (0,2.8) {$\VDD$};
\draw (pa.D) -- (-1.3,1.2) -- (1.3,1.2) -- (pb.D);\draw (0,1.2) -- (pc.S);
\draw (pc.D) -- (0,-.4) -- (2.2,-.4) node[right]{$Y$};
\node[nmos] (na) at (-1.3,-1.3) {};\node[nmos] (nb) at (-1.3,-2.9) {};\node[nmos] (nc) at (1.3,-2.1) {};
\draw (na.D) -- (-1.3,-.4) -- (1.3,-.4) -- (nc.D);\draw (na.S) -- (nb.D);
\draw (nb.S) -- (-1.3,-3.75) -- (1.3,-3.75) -- (nc.S);\draw (0,-3.75) node[ground]{};
\draw (pa.G) node[left]{$A$};\draw (pb.G) node[left]{$B$};\draw (pc.G) node[left]{$C$};
\draw (na.G) node[left]{$A$};\draw (nb.G) node[left]{$B$};\draw (nc.G) node[left]{$C$};
\end{circuitikz}}
% Figuras vetoriais autorais: parâmetros controlam etapas da apresentação.
\definecolor{eletron}{RGB}{0,83,134}
\definecolor{lacuna}{RGB}{199,123,0}
\definecolor{ativo}{RGB}{0,132,75}
\tikzset{fio/.style={line width=.9pt,black!65},ligado/.style={line width=2pt,ativo},cortado/.style={line width=.7pt,black!25},anotacao/.style={font=\small,align=center},passo/.style={font=\scriptsize,text=blue_theme}}
\newcommand{\chamada}[1]{\par\medskip\centering\textcolor{blue_theme}{\textbf{#1}}\par}
\newcommand{\etapa}[2]{\par\smallskip{\scriptsize\color{blue_theme}Etapa #1\enspace•\enspace #2}\par}
% PN: 0 contato inicial; 1 difusão; 2 equilíbrio; 3 direta; 4 reversa.
\newcommand{\juncaoPN}[1]{%
\begin{tikzpicture}[x=.87cm,y=.85cm,font=\small]
\path[use as bounding box] (-.3,-1.65) rectangle (10.3,4.9);
\fill[lacuna!9] (0,0) rectangle (5,3);\fill[eletron!9] (5,0) rectangle (10,3);
\node[blue_theme] at (2,3.5){\textbf{Região P}};\node[blue_theme] at (8,3.5){\textbf{Região N}};
\ifnum#1>1
 \def\depl{1.25}\ifnum#1=3\def\depl{.6}\fi\ifnum#1=4\def\depl{2}\fi
 \fill[gray!25] (5-\depl,0) rectangle (5+\depl,3);
 \draw[densely dashed,gray] (5-\depl,0)--(5-\depl,3);\draw[densely dashed,gray] (5+\depl,0)--(5+\depl,3);
 \foreach \y in {.55,1.25,1.95,2.6}{\node[eletron] at (5-.5*\depl,\y){$-$};\node[lacuna] at (5+.5*\depl,\y){$+$};}
 \draw[<->,blue_theme] (5-\depl,4.15)--node[above]{\strut Depleção}(5+\depl,4.15);
 \draw[-{Latex},very thick,blue_theme] (5+\depl,-.55)--node[below]{Campo interno}(5-\depl,-.55);
\fi
\foreach \x/\y in {.7/.6,1.6/1.4,2.5/.65,.8/2.3,2.6/2.4}{\draw[lacuna,fill=white,line width=1pt] (\x,\y) circle(.13);}
\foreach \x/\y in {7.5/.6,8.4/1.4,9.3/.65,7.4/2.4,9.2/2.3}{\fill[eletron] (\x,\y) circle(.12);}
\ifnum#1<2
 \foreach \x/\y in {3.6/.8,4.3/2,3.4/2.5}{\draw[lacuna,fill=white,line width=1pt] (\x,\y) circle(.13);}
 \foreach \x/\y in {5.7/.8,6.4/2,6.6/2.5}{\fill[eletron] (\x,\y) circle(.12);}
\fi
\ifnum#1=1
 \draw[-{Latex},very thick,eletron] (6.3,1.2)--node[above]{elétrons}(3.8,1.2);
 \draw[-{Latex},very thick,lacuna] (3.6,2.3)--node[above]{lacunas}(6.1,2.3);
 \node[font=\scriptsize] at (5,-1.25){Difusão e recombinação próximas à interface};
\fi
\ifnum#1=0\node at (5,-1){\textcolor{lacuna}{$\circ$} lacuna \qquad\textcolor{eletron}{$\bullet$} elétron};\fi
\ifnum#1=2\node[font=\scriptsize] at (5,-1.3){Sinais na faixa cinza: íons fixos, não portadores móveis};\fi
\ifnum#1=3
 \node at (1.2,4.3){$+$};\node at (8.8,4.3){$-$};
 \node[font=\scriptsize] at (5,-1.3){Direta: barreira menor, corrente maior};
\fi
\ifnum#1=4
 \node at (1.2,4.3){$-$};\node at (8.8,4.3){$+$};
 \node[font=\scriptsize] at (5,-1.3){Reversa: barreira maior, pequena fuga antes da ruptura};
\fi
\end{tikzpicture}}
% NMOS planar: 0 sem canal forte; 1 depleção; 2 canal; 3 corrente.
\newcommand{\moscanal}[1]{%
\begin{tikzpicture}[x=.85cm,y=.8cm,font=\small]
\path[use as bounding box] (-.3,-1.1) rectangle (11,5.2);
\fill[eletron!8,draw=blue_theme] (0,0) rectangle (10,2.8);
\fill[yellow!35,draw=blue_theme] (.5,1.7) rectangle (2.5,2.8);
\fill[yellow!35,draw=blue_theme] (7.5,1.7) rectangle (9.5,2.8);
\node at (1.5,2.2){$n^+$};\node at (8.5,2.2){$n^+$};
\fill[gray!40,draw=blue_theme] (2.7,2.85) rectangle (7.3,3.08);
\fill[eletron!55,draw=blue_theme] (2.7,3.08) rectangle (7.3,3.55);
\node at (5,.5){Corpo P};\node at (1.5,3.9){Fonte};\node at (8.5,3.9){Dreno};\node at (5,4){Porta};
\draw (1.5,2.8)--(1.5,3.6);\draw (8.5,2.8)--(8.5,3.6);\draw (5,3.55)--(5,3.8);
\node[font=\scriptsize,anchor=west] at (8.1,4.72){Dielétrico};\draw[-{Latex},black!50] (8.1,4.35)--(7.65,4.35)--(7.25,2.97);
\ifnum#1<1
 \foreach \x/\y in {3.3/1.2,4.3/2.1,5.2/1.4,6.5/2.1}{\draw[lacuna,fill=white] (\x,\y) circle(.12);}
 \node at (5,4.75){$V_{GS}$ baixo};
\else
 \fill[gray!20] (2.5,1.55) rectangle (7.5,2.8);
 \foreach \x in {3,4,5,6,7}{\node[eletron,font=\scriptsize] at (\x,2.1){$-$};}
 \foreach \x in {3.2,4.4,5.6,6.8}{\node[red] at (\x,3.32){$+$};}
 \ifnum#1=1
  \node at (5,4.75){A tensão de porta aumenta};
  \draw[-{Latex},lacuna,thick] (4,1.55)--(4,.8);\draw[-{Latex},lacuna,thick] (6,1.55)--(6,.8);
  \node[font=\scriptsize] at (5,-.65){Lacunas se afastam da superfície: depleção};
 \else
  \fill[eletron!70] (2.5,2.57) rectangle (7.5,2.8);
  \foreach \x in {2.9,3.6,4.3,5,5.7,6.4,7.1}{\node[white,font=\scriptsize] at (\x,2.69){$-$};}
  \node at (5,4.75){$V_{GS}>V_{TN}$: canal de elétrons};
  \ifnum#1=3
   \draw[-{Latex},white,line width=3.5pt] (8.4,2.69)--(1.6,2.69);\draw[-{Latex},line width=2pt,ativo] (8.4,2.69)--(1.6,2.69);\node[ativo] at (5,1.15){Corrente convencional $I_D$};
   \node[font=\scriptsize] at (5,-.65){Com $V_D>V_S$, há corrente pelo canal};
  \else\node[font=\scriptsize] at (5,-.65){Inversão: a superfície conduz por elétrons};\fi
 \fi
\fi
\end{tikzpicture}}
% Circuitos por estado: caminho ativo verde; transistor bloqueado cinza.
\newcommand{\inversorestado}[1]{%
\begin{circuitikz}[scale=.9,transform shape]
\path[use as bounding box] (-3,-2.9) rectangle (3,3);
\ifnum#1=0\def\pstyle{ativo}\def\nstyle{black!25}\def\saida{1}\else\def\pstyle{black!25}\def\nstyle{ativo}\def\saida{0}\fi
\node[pmos,color=\pstyle] (p) at (0,1.3) {};\node[nmos,color=\nstyle] (n) at (0,-1.1) {};
\draw[ligado,color=\pstyle] (p.S)--(0,2.2) node[above,black]{$\VDD$};
\draw[ligado,color=\pstyle] (p.D)--(0,.1);
\draw[ligado,color=\nstyle] (n.D)--(0,.1);\draw[ligado,color=\nstyle] (n.S)--(0,-2.1) node[ground,black]{};
\draw[fio] (p.G)--(-1.5,1.3)--(-1.5,-1.1)--(n.G);\draw[fio] (-1.5,.1)--(-2.1,.1) node[left]{$A=#1$};
\draw[fio] (0,.1) to[short,*-] (1.3,.1) node[right]{$Y=\saida$};
\node[font=\small,text=\pstyle,right] at (.5,1.3){PMOS};\node[font=\small,text=\nstyle,right] at (.5,-1.1){NMOS};
\node[font=\scriptsize,ativo] at (0,-2.8){Verde: caminho que fixa a saída};
\end{circuitikz}}
\newcommand{\portaestado}[3]{%
\begin{circuitikz}[scale=.78,transform shape]
\path[use as bounding box] (-3,-3.8) rectangle (3.7,3.6);
\ifnum#2=0\def\pa{ativo}\def\na{black!25}\else\def\pa{black!25}\def\na{ativo}\fi
\ifnum#3=0\def\pb{ativo}\def\nb{black!25}\else\def\pb{black!25}\def\nb{ativo}\fi
\ifnum#1=0
 \pgfmathtruncatemacro{\out}{1-#2*#3}
 \node[pmos,color=\pa] (pa) at (-1.2,1.7) {};\node[pmos,color=\pb] (pb) at (1.4,1.7) {};
 \node[nmos,color=\na] (na) at (0,-.7) {};\node[nmos,color=\nb] (nb) at (0,-2.3) {};
 \draw[fio] (pa.S)--(-1.2,2.6)--(1.4,2.6)--(pb.S);\node[above] at (.1,2.6){$\VDD$};
 \draw[ligado,color=\pa] (pa.D)--(-1.2,.5)--(0,.5);\draw[ligado,color=\pb] (pb.D)--(1.4,.5)--(0,.5);
 \draw[ligado,color=\na] (0,.5)--(na.D);\draw[ligado,color=\na] (na.S)--(0,-1.5);
 \draw[ligado,color=\nb] (0,-1.5)--(nb.D);\draw[ligado,color=\nb] (nb.S)--(0,-3.2) node[ground,black]{};
 \draw[fio] (0,.5) to[short,*-] (2.2,.5) node[right]{$Y=\out$};
\else
 \pgfmathtruncatemacro{\out}{(1-#2)*(1-#3)}
 \node[pmos,color=\pa] (pa) at (0,2) {};\node[pmos,color=\pb] (pb) at (0,.4) {};
 \node[nmos,color=\na] (na) at (-1.2,-1.8) {};\node[nmos,color=\nb] (nb) at (1.4,-1.8) {};
 \draw[ligado,color=\pa] (pa.S)--(0,2.85) node[above,black]{$\VDD$};\draw[ligado,color=\pa] (pa.D)--(0,1.2);
 \draw[ligado,color=\pb] (0,1.2)--(pb.S);\draw[ligado,color=\pb] (pb.D)--(0,-.6);
 \draw[ligado,color=\na] (na.D)--(-1.2,-.6)--(0,-.6);\draw[ligado,color=\nb] (nb.D)--(1.4,-.6)--(0,-.6);
 \draw[ligado,color=\na] (na.S)--(-1.2,-2.8);\draw[ligado,color=\nb] (nb.S)--(1.4,-2.8);
 \draw[fio] (-1.2,-2.8)--(1.4,-2.8);\draw (0,-2.8) node[ground]{};
 \draw[fio] (0,-.6) to[short,*-] (2.2,-.6) node[right]{$Y=\out$};
\fi
\draw (pa.G) node[left,black]{$A=#2$};\draw (pb.G) node[left,black]{$B=#3$};
\draw (na.G) node[left,black]{$A=#2$};\draw (nb.G) node[left,black]{$B=#3$};
\end{circuitikz}}

% Figuras originais para a abertura histórica. Referências: dados/FONTES_CONTEXTO.md.
% As cores distinguem Gate, canal e interconexões; desenhos sem escala.
\newcommand{\linhaHistorica}{%
\begin{tikzpicture}[x=1cm,y=1cm,font=\small]
\draw[seta] (0,0)--(10.3,0);
\foreach \x/\ano/\titulo/\sub in {0.4/1947/Transistor/Bell Labs,2.8/1958--59/Circuito integrado/Kilby e Noyce,5.2/1960--63/MOS e CMOS/Controle por campo,7.6/1971/Microprocessador/Intel 4004,10/Hoje/Sistemas integrados/Múltiplas funções}{%
\fill[blue_theme] (\x,0) circle(.09);
\node[blue_theme,above=3mm,align=center,font=\small\bfseries] at (\x,0){\ano};
\node[below=3mm,align=center,text width=2.25cm,font=\scriptsize] at (\x,0){\textbf{\titulo}\\[1mm]\sub};}
% Dispositivo discreto, circuito, CPU, SoC: motivos conceituais.
\draw[blue_theme,thick] (.05,1.5) rectangle (.75,2.25);
\foreach \x in {.15,.4,.65}{\draw[blue_theme] (\x,1.5)--(\x,1.1);}
\draw[blue_theme,thick] (2,1.2) rectangle (3.6,2.5);
\foreach \x in {2.2,2.7,3.2}{\draw[blue_theme] (\x,1.6)--++(.2,.45)--++(.2,-.45);}
\draw[blue_theme,thick] (4.6,1.4)--(5.8,1.4);\draw[blue_theme,thick] (4.6,1.65)--(5.8,1.65);\draw[blue_theme,line width=4pt] (4.9,2)--(5.5,2);
\draw[blue_theme,thick] (6.8,1.2) rectangle (8.4,2.5);\node at (7.6,1.85){CPU};
\draw[blue_theme,thick] (9.1,1.2) rectangle (10.9,2.5);
\foreach \x/\y in {9.25/1.35,10.05/1.35,9.25/1.95,10.05/1.95}{\fill[blue_theme!20] (\x,\y) rectangle ++(.65,.4);}
\path[use as bounding box] (-.8,-1.3) rectangle (11.2,2.9);
\end{tikzpicture}}

\newcommand{\graficoMoore}{%
\begin{tikzpicture}[x=.86cm,y=.19cm,font=\small]
\draw[->,blue_theme] (0,0)--(9,0) node[right]{anos};
\draw[->,blue_theme] (0,0)--(0,19) node[above,align=center]{quantidade\\relativa};
\foreach \y in {1,2,4,8,16}{\draw[gray!25] (0,\y)--(8,\y);\node[left] at (0,\y){\y};}
\foreach \x in {0,2,4,6,8}{\draw (\x,0)--(\x,-.4) node[below]{\x};}
\draw[blue_theme,very thick,domain=0:8,samples=60] plot (\x,{2^(\x/2)});
\foreach \x/\y in {0/1,2/2,4/4,6/8,8/16}{\fill[blue_theme] (\x,\y) circle[x radius=.07,y radius=.32];}
\node[align=center,fill=white,inner sep=2mm] at (3.6,14){\textbf{Dobrar a cada dois anos}\\$N(t)=N_0\,2^{t/2}$};
\path[use as bounding box] (-1.4,-3) rectangle (10.2,24);
\end{tikzpicture}}

\newcommand{\graficoProcessos}{%
\begin{tikzpicture}[x=.175cm,y=.79cm,font=\scriptsize]
\fill[blue_theme!6] (47,-.1) rectangle (56,4.8);
\draw[->,blue_theme] (0,0)--(57,0);\draw[->,blue_theme] (0,0)--(0,4.8);
\node[anchor=west] at (0,5.15){Referência de processo (nm) · eixo logarítmico};
\foreach \y/\t in {0/1,1/10,2/100,3/{1\,000},4/{10\,000}}{\draw[gray!25] (0,\y)--(56,\y);\node[left] at (0,\y){\t};}
\draw[blue_theme,very thick] (0,4)--(7,3.4771)--(14,3.1761)--(22,2.9031)--(29,2.2553)--(37,1.6532);
\foreach \x/\y/\valor/\ano in {0/4/{10\,000}/1971,14/3.1761/{1\,500}/1985,29/2.2553/180/2000,37/1.6532/45/2008}{%
\fill[blue_theme] (\x,\y) circle[x radius=.45,y radius=.1];
\ifnum\x>0\node[above=1.5mm] at (\x,\y){\valor};\fi\node[below=2mm] at (\x,0){\ano};}
\foreach \x/\y/\nome/\ano in {49/.699/N5/2020,54/.301/N2/2025}{%
\draw[lacuna,thick,fill=white] (\x,\y) circle[x radius=.5,y radius=.11];
\node[above=2mm,lacuna] at (\x,\y){\nome};\node[below=2mm] at (\x,0){\ano};}
\node[align=center,text width=1.7cm] at (51.5,3.4){Nós comerciais\\recentes};
\node[anchor=west] at (0,-1.1){$1\,\mu$m = $1\,000$ nm};
\path[use as bounding box] (-7,-1.5) rectangle (58,5.6);
\end{tikzpicture}}

\newcommand{\litografiaEtapa}[1]{%
\begin{tikzpicture}[x=1cm,y=1cm,font=\small]
\fill[blue_theme!15] (0,0) rectangle (7,.7);\node at (3.5,.32){substrato};
\ifnum#1<3
\fill[blue_theme!65] (0,.7) rectangle (7,.95);
\else
\foreach \x in {0,2.8,5.6}{\fill[blue_theme!65] (\x,.7) rectangle ++(1.4,.25);}
\fi
\ifnum#1=1
\fill[yellow!70] (0,.95) rectangle (7,1.3);
\foreach \x in {1.4,4.2}{\fill[red!12] (\x,.95) rectangle ++(1.4,.35);\fill[red!10] (\x,1.4) rectangle ++(1.4,2.2);\draw[-{Latex},red] (\x+.7,3.5)--(\x+.7,1.45);}
\draw[blue_theme,line width=4pt] (0,3.8)--(1.4,3.8) (2.8,3.8)--(4.2,3.8) (5.6,3.8)--(7,3.8);
\node[left] at (0,3.8){máscara};\node[right,align=left,font=\scriptsize] at (7.2,2.6){Luz projetada\\sobre o resiste};
\else
\foreach \x in {0,2.8,5.6}{\fill[yellow!70] (\x,.95) rectangle ++(1.4,.35);}
\ifnum#1=2
\foreach \x in {2.1,4.9}{\draw[seta] (\x,2.6)--(\x,1.5);}
\node[align=center] at (3.5,3.5){Revelação remove as partes expostas\\de um resiste positivo.};
\else
\foreach \x in {2.1,4.9}{\draw[-{Latex},thick,red] (\x,2.6)--(\x,.85);}
\node[align=center] at (3.5,3.5){Gravação transfere o padrão\\à camada abaixo do resiste.};
\fi\fi
\draw[blue_theme] (0,0) rectangle (7,.7);
\draw[-,blue_theme] (7,.82)--(7.6,.82) node[right,font=\scriptsize]{camada a padronizar};
\draw[-,lacuna] (6.4,1.2)--(7.6,1.7) node[right,font=\scriptsize]{resiste};
\path[use as bounding box] (-1.5,-.4) rectangle (10.1,4.5);
\end{tikzpicture}}

\newcommand{\transistoresModernos}{%
\begin{tikzpicture}[x=1cm,y=1cm,font=\small]
\foreach \x/\nome in {0/Planar,3.8/FinFET,7.6/{GAA · nanosheets}}{%
\node[blue_theme,font=\small\bfseries] at (\x+1.25,3.2){\nome};
\fill[gray!12] (\x,0) rectangle ++(2.5,.5);}
% Cortes transversais ao canal. Dielétrico branco entre Gate e silício.
\fill[ativo!60] (.4,.5) rectangle (2.1,.85);
\fill[blue_theme!70] (.35,1) rectangle (2.15,1.5);
\fill[blue_theme!70] (4.35,.55) rectangle (5.75,2.4);
\fill[white] (4.75,.5) rectangle (5.35,2.05);
\fill[ativo!60] (4.87,.5) rectangle (5.23,1.93);
\fill[blue_theme!70] (8.1,.65) rectangle (9.6,2.55);
\foreach \y in {.9,1.5,2.1}{\fill[white,rounded corners=2pt] (8.3,\y) rectangle ++(1.1,.3);\fill[ativo!60,rounded corners=1pt] (8.4,\y+.08) rectangle ++(.9,.14);}
\node[align=center,font=\scriptsize] at (1.25,-.5){Gate sobre\\a superfície};
\node[align=center,font=\scriptsize] at (5.05,-.5){Gate envolve\\três faces};
\node[align=center,font=\scriptsize] at (8.85,-.5){Gate envolve\\cada canal};
\fill[blue_theme!70] (2,-1.45) rectangle ++(.25,.2);\node[right,font=\scriptsize] at (2.3,-1.35){Gate};
\fill[ativo!60] (5,-1.45) rectangle ++(.25,.2);\node[right,font=\scriptsize] at (5.3,-1.35){canal};
\path[use as bounding box] (-.3,-1.7) rectangle (10.4,3.6);
\end{tikzpicture}}

\newcommand{\pilhaCMOSdois}{%
\begin{tikzpicture}[x=1cm,y=1cm,font=\small]
\foreach \y/\cor/\titulo in {0/gray/Alimentação e acesso externo,1.2/ativo/Lógica densa,2.4/blue_theme/Memória,3.6/lacuna/Interfaces e funções especializadas}{%
\fill[\cor!14] (0,\y)--(5,\y)--(5.8,\y+.45)--(.8,\y+.45)--cycle;
\draw[\cor,thick] (0,\y)--(5,\y)--(5.8,\y+.45)--(.8,\y+.45)--cycle;
\node[anchor=west,font=\scriptsize] at (6,\y+.2){\titulo};}
\foreach \x in {1,2,3,4}{\draw[blue_theme,thick,densely dashed] (\x,.5)--(\x,3.5);}
\node[anchor=west,font=\scriptsize,align=left] at (6,-.9){Conexões 3D densas\\entre camadas funcionais};
\path[use as bounding box] (-.3,-1.4) rectangle (10.6,4.5);
\end{tikzpicture}}

% Diagramas didáticos originais. Sem parâmetros de um PDK.
\newcommand{\corteCMOS}[1]{%
\begin{tikzpicture}[x=.88cm,y=.88cm,font=\small]
\path[use as bounding box] (-.4,-.7) rectangle (11.4,5.4);
\fill[blue_theme!10] (0,0) rectangle (11,1.4);
\fill[lacuna!17] (.55,.2) rectangle (5.05,1.4);
\draw[lacuna,thick] (.55,1.4)--(.55,.2)--(5.05,.2)--(5.05,1.4);
\fill[gray!35] (5.35,.9) rectangle (6.45,1.4);
\node[font=\scriptsize] at (2.8,.5){poço N};
\node[font=\scriptsize] at (8.1,.45){substrato P};
\node[blue_theme,font=\bfseries] at (2.8,-.35){PMOS};
\node[blue_theme,font=\bfseries] at (8.55,-.35){NMOS};
\ifnum#1>1
\fill[white] (0,1.4) rectangle (11,1.54);
\foreach \x/\tipo/\cor in {1.1/{p$^+$}/lacuna,3.5/{p$^+$}/lacuna,7/{n$^+$}/eletron,9.4/{n$^+$}/eletron}{%
\fill[\cor!40] (\x,1) rectangle ++(.9,.4);\node[font=\scriptsize] at (\x+.45,1.18){\tipo};}
\foreach \x in {2.2,8.15}{\fill[lacuna!65] (\x,1.43) rectangle ++(.9,.1);\fill[blue_theme!65] (\x,1.59) rectangle ++(.9,.38);}
\ifnum#1=2
\draw[blue_theme] (2.65,2.05)--(2.65,2.6) node[above,font=\scriptsize]{porta};
\draw[lacuna] (8.6,1.48)--(9.5,2.5) node[above,font=\scriptsize]{dielétrico};
\node[font=\scriptsize,align=center] at (5.5,3.9){Os dois dispositivos têm regiões de fonte e dreno\\e uma porta isolada do semicondutor.};
\fi
\else
\node[font=\scriptsize,align=center] at (5.5,3.5){Regiões com dopagens diferentes\\permitem construir os dispositivos complementares.};
\draw[blue_theme] (5.9,1.4)--(5.9,2.15) node[above,font=\scriptsize]{isolação};
\fi
\ifnum#1=3
\foreach \x in {1.55,3.95,7.45,9.85}{\fill[black!55] (\x-.075,1.4) rectangle (\x+.075,2.55);}
\draw[line width=2pt,ativo] (1.55,2.55)--(1.55,3.55) node[above]{$\VDD$};
\draw[line width=2pt,ativo] (9.85,2.55)--(9.85,3.55) node[above]{GND};
\draw[line width=2pt,blue_theme] (3.95,2.55)--(3.95,3)--(7.45,3)--(7.45,2.55);
\draw[line width=2pt,blue_theme] (5.8,3)--(5.8,3.55) node[above]{$Y$};
\draw[line width=1.6pt,lacuna] (2.65,1.97)--(2.65,4.7)--(8.6,4.7)--(8.6,1.97);
\node[above,lacuna] at (5.6,4.7){entrada $A$};
\node[font=\scriptsize,anchor=west] at (0,2.55){\strut};
\fi
\draw[blue_theme] (0,0) rectangle (11,1.4);
\end{tikzpicture}}

\newcommand{\cicloFabricacao}{%
\begin{tikzpicture}[x=1cm,y=1cm,font=\small]
\foreach \x in {0,3.2,6.4}{\fill[blue_theme!15] (\x,1.5) rectangle ++(2.2,.55);}
\fill[blue_theme!65] (0,2.05) rectangle (2.2,2.4);
\fill[blue_theme!65] (3.2,2.05) rectangle (5.4,2.4);
\foreach \x in {3.2,4.65}{\fill[lacuna!60] (\x,2.4) rectangle ++(.75,.3);}
\foreach \x in {6.4,7.85}{\fill[blue_theme!65] (\x,2.05) rectangle ++(.75,.35);}
\node[font=\bfseries,blue_theme] at (1.1,3.45){Criar uma camada};
\node[font=\bfseries,blue_theme] at (4.3,3.45){Definir um padrão};
\node[font=\bfseries,blue_theme] at (7.5,3.45){Modificar o material};
\node[align=center,font=\scriptsize] at (1.1,.7){Crescer ou depositar\\isolantes e condutores};
\node[align=center,font=\scriptsize] at (4.3,.7){Litografia no resiste\\e alinhamento};
\node[align=center,font=\scriptsize] at (7.5,.7){Gravar, dopar ou abrir\\acesso às camadas};
\draw[seta] (2.4,2.15)--(3,2.15);\draw[seta] (5.6,2.15)--(6.2,2.15);
\draw[seta] (7.5,.1)--(7.5,-.6)--(1.1,-.6)--(1.1,.1);
\node[fill=white,inner sep=2pt,font=\scriptsize] at (4.3,-.6){Repetir com novas máscaras e materiais};
\path[use as bounding box] (-.5,-1) rectangle (9,4);
\end{tikzpicture}}

\newcommand{\escalaAtomica}{%
\begin{tikzpicture}[x=1cm,y=1cm,font=\small]
\node[blue_theme,font=\bfseries] at (1.9,4.1){Tamanho de referência};
\shade[inner color=blue_theme!5,outer color=blue_theme!35] (1.9,1.9) circle(.85);
\node[blue_theme,font=\large] at (1.9,1.9){Si};
\draw[<->,blue_theme] (1.05,.65)--node[below,font=\scriptsize]{diâmetro covalente $\approx0{,}228$ nm}(2.75,.65);
\node[blue_theme,font=\bfseries] at (7.2,4.1){Distância na rede cristalina};
\draw[blue_theme!40,line width=2pt] (6,1.9)--(8.35,1.9);
\foreach \x in {6,8.35}{\shade[inner color=blue_theme!5,outer color=blue_theme!35] (\x,1.9) circle(1.14);\node[blue_theme] at (\x,1.9){Si};}
\draw[densely dashed,blue_theme!50] (6,1.9)--(6,3.45);\draw[densely dashed,blue_theme!50] (8.35,1.9)--(8.35,3.45);
\draw[<->,blue_theme] (6,3.35)--node[above,font=\scriptsize]{$\approx0{,}235$ nm}(8.35,3.35);
\node[align=center,font=\scriptsize] at (7.2,.3){Separação entre vizinhos mais próximos.\\A organização real é tridimensional.};
\path[use as bounding box] (-.2,-.4) rectangle (10,4.6);
\end{tikzpicture}}

\newcommand{\reguaAtomica}{%
\begin{tikzpicture}[x=.74cm,y=1cm,font=\small]
\node[anchor=west,blue_theme,font=\bfseries] at (0,4.3){Comprimentos físicos em uma mesma escala};
\draw[<->,lacuna,thick] (0,3.5)--node[above,font=\scriptsize]{13,5 nm · comprimento de onda EUV}(13.5,3.5);
\draw[<->,blue_theme,thick] (0,2.5)--node[above,font=\scriptsize]{8 nm · resolução nominal High-NA}(8,2.5);
\draw[blue_theme!35,line width=.8pt] (0,1.5)--(8,1.5);
\foreach \i in {0,...,34}{\fill[blue_theme!65] ({\i*.235},1.5) circle[x radius=.095,y radius=.07];}
\node[anchor=west,font=\scriptsize] at (0,.9){8 nm $\approx$ 34 separações de 0,235 nm};
\node[anchor=west,align=left,font=\scriptsize,text width=3.1cm] at (9.1,1.9){\textbf{``2 nm'' é o nome do nó.}\\Não é uma medida física única do transistor.};
\draw[blue_theme!30] (8.65,.7)--(8.65,2.8);
\node[anchor=west,font=\scriptsize] at (0,-.05){Régua unidimensional ilustrativa; não representa a estrutura de um transistor.};
\path[use as bounding box] (-.4,-.4) rectangle (14,4.7);
\end{tikzpicture}}

\newcommand{\barreiraTunelamento}[1]{%
\begin{tikzpicture}[x=1cm,y=1cm,font=\small]
\def\largura{2.7}\ifnum#1=2\def\largura{.85}\fi
\draw[-{Latex},blue_theme] (.2,.3)--(.2,4) node[above,font=\scriptsize]{energia};
\fill[blue_theme!12] (3.4,.55) rectangle (3.4+\largura,3.2);
\draw[blue_theme,very thick] (.5,.55)--(3.4,.55)--(3.4,3.2)--(3.4+\largura,3.2)--(3.4+\largura,.55)--(10.1,.55);
\draw[lacuna,dashed] (.5,1.6)--(10.1,1.6);
\node[anchor=west,font=\scriptsize,lacuna] at (.7,1.95){energia do elétron};
\draw[-{Latex},eletron,very thick] (.8,1.6)--(2.9,1.6);
\fill[eletron] (1.1,1.6) circle(.09);
\draw[<->,blue_theme] (3.4,3.65)--node[above,font=\scriptsize]{espessura da barreira}(3.4+\largura,3.65);
\ifnum#1=1
\draw[-{Latex},eletron!35,dashed,line width=.6pt] (6.3,1.6)--(7.5,1.6);
\node[font=\scriptsize,align=center] at (8,2.65){Transmissão menor\\para esta barreira};
\else
\draw[-{Latex},red,very thick] (4.6,1.6)--(8,1.6);
\node[font=\scriptsize,align=center] at (7.3,2.65){Maior probabilidade\\de atravessar};
\fi
\node[font=\scriptsize,anchor=west] at (.7,-.1){Diagrama de energia: a altura representa energia, não uma parede física.};
\path[use as bounding box] (-.2,-.55) rectangle (10.5,4.35);
\end{tikzpicture}}

% Analogia urbana original: ilustra conectividade; não fixa custos ou ganhos.
\newcommand{\arvoreUrbana}[2]{%
\draw[brown!70!black,line width=1pt] (#1,#2)--++(0,.4);
\fill[green_ufc!45!black] (#1,#2+.5) circle(.17);
\fill[green_ufc!65] (#1-.1,#2+.48) circle(.14);}
\newcommand{\predioUrbano}[5]{%
\begin{scope}[shift={(#1,#2)}]
\fill[black!8] (.12,-.08)--(#3+.52,.1)--(#3+.7,.35)--(.4,.2)--cycle;
\fill[#5!18] (0,0) rectangle (#3,#4);
\fill[#5!38] (#3,0)--(#3+.35,.25)--(#3+.35,#4+.25)--(#3,#4)--cycle;
\fill[#5!8] (0,#4)--(.35,#4+.25)--(#3+.35,#4+.25)--(#3,#4)--cycle;
\draw[#5!65!black,line width=.5pt] (0,0)--(0,#4)--(.35,#4+.25)--(#3+.35,#4+.25)--(#3+.35,.25)--(#3,0)--cycle;
\draw[#5!60!black] (0,#4)--(#3,#4)--(#3,0) (#3,#4)--(#3+.35,#4+.25);
\foreach \yy in {.45,.9,1.35}{\ifdim\yy pt<#4pt\relax\foreach \xx in {.2,.6}{\ifdim\xx pt<#3pt\relax\fill[white!90!#5] (\xx,\yy) rectangle ++(.2,.22);\fi}\fi}
\fill[#5!60!black] (.36,0) rectangle (.61,.34);
\end{scope}}
\newcommand{\pessoaDados}[2]{%
\begin{scope}[shift={(#1,#2)}]
\fill[blue_theme] (0,.27) circle(.065);
\draw[blue_theme,line width=1pt] (0,.2)--(0,.03)--(-.1,-.13) (0,.03)--(.1,-.13) (-.12,.13)--(0,.17)--(.17,.14);
\fill[lacuna!75] (.13,.02) rectangle (.33,.22);\draw[lacuna!70!black] (.13,.02) rectangle (.33,.22);
\end{scope}}
\newcommand{\cidadeDados}[1]{%
\begin{tikzpicture}[x=1cm,y=1cm,font=\small,line join=round]
\ifnum#1=1
% Terreno, calçadas e praças delimitam bairros, sem cartões de texto.
\fill[green_ufc!5] (-.1,.45) rectangle (11.7,4.45);
\fill[blue_theme!6] (.15,1.4) rectangle (3.65,4.2);
\fill[ativo!8] (7.75,1.4) rectangle (11.4,4.2);
\fill[gray!15] (3.8,2.5) rectangle (7.5,4.1);
\foreach \x/\y in {4.1/2.8,5/2.8,5.9/2.8,6.8/2.8,4.1/3.6,6.8/3.6}{\arvoreUrbana{\x}{\y}}
% Via sinuosa: conserva o significado de distância e cruzamentos.
\draw[gray!55,line width=14pt] (2.1,1.8)--(3.5,1.8)--(3.5,.9)--(7.9,.9)--(7.9,1.8)--(9.3,1.8);
\draw[white,line width=.65pt,dash pattern=on 4pt off 3pt] (2.1,1.8)--(3.5,1.8)--(3.5,.9)--(7.9,.9)--(7.9,1.8)--(9.3,1.8);
\foreach \x in {3.32,3.5,3.68}{\draw[white,line width=1.4pt] (\x,1.56)--(\x,2.04);}
\foreach \x in {7.72,7.9,8.08}{\draw[white,line width=1.4pt] (\x,1.56)--(\x,2.04);}
% Edifícios e depósitos do bairro de memória.
\predioUrbano{.55}{2.15}{1.15}{1.8}{blue_theme}
\predioUrbano{2.1}{2.15}{.95}{1.35}{blue_theme}
\predioUrbano{8.45}{2.15}{1.15}{1.8}{ativo}
\predioUrbano{10}{2.15}{.95}{1.35}{ativo}
\foreach \x/\y in {.35/1.8,3.2/3.7,8.25/3.7,11.1/1.8}{\arvoreUrbana{\x}{\y}}
% Caminhão de dados com cabine, carga e rodas.
\begin{scope}[shift={(5.2,.96)}]
\fill[blue_theme] (-.55,-.02) rectangle (.3,.43);\fill[blue_theme!70] (.3,-.02)--(.7,-.02)--(.7,.24)--(.52,.43)--(.3,.43)--cycle;
\fill[white] (.38,.24) rectangle (.57,.34);
\foreach \x in {-.35,.48}{\fill[black!75] (\x,-.08) circle(.1);\fill[white] (\x,-.08) circle(.035);}
\node[font=\tiny,text=white] at (-.12,.2){0101};
\end{scope}
\draw[-{Latex},blue_theme,very thick] (4.1,.43)--(7.2,.43);
\node[font=\scriptsize,blue_theme,fill=white,inner sep=2pt] at (5.7,.05){Dados atravessam a cidade};
\node[blue_theme,font=\bfseries] at (1.9,4.9){Bairro 1 · memória};
\node[ativo,font=\bfseries] at (9.6,4.9){Bairro 2 · processamento};
\node[font=\scriptsize,black!65] at (5.7,4.45){Mais percurso e interfaces de comunicação};
\else
% Corte de um prédio real: fachada lateral, pilares, lajes e três pavimentos.
\fill[black!8] (2.5,-.15)--(9.2,-.15)--(10.3,.3)--(3.6,.3)--cycle;
\fill[blue_theme!18] (8.55,.2)--(9.2,.65)--(9.2,5.1)--(8.55,4.65)--cycle;
\foreach \y in {.5,1.8,3.1}{\foreach \x in {8.67,8.92}{\fill[white!85!blue_theme] (\x,\y+.15)--(\x+.13,\y+.24)--(\x+.13,\y+.73)--(\x,\y+.64)--cycle;}}
\fill[blue_theme!6] (2.8,.25) rectangle (8.5,4.65);
\fill[ativo!8] (2.85,1.65) rectangle (7.5,3.05);
\fill[blue_theme!10] (2.85,3.1) rectangle (7.5,4.5);
\foreach \y in {.2,1.6,3.05,4.5}{%
\fill[gray!45] (2.65,\y) rectangle (8.55,\y+.13);
\fill[gray!22] (2.65,\y+.13)--(3.3,\y+.55)--(9.2,\y+.55)--(8.55,\y+.13)--cycle;}
% Fundo dos andares desenhado à frente apenas com os elementos internos.
\foreach \x in {2.7,7.5,8.4}{\fill[gray!45] (\x,.2) rectangle (\x+.11,4.65);}
% Memória: três estantes com módulos, não janelas abstratas.
\foreach \x in {3.15,4.5,5.85}{%
\fill[blue_theme!60] (\x,3.28) rectangle ++(.9,.95);
\foreach \y in {3.4,3.66,3.92}{\fill[white!85!blue_theme] (\x+.12,\y) rectangle ++(.66,.13);\fill[lacuna] (\x+.65,\y+.035) circle(.025);}}
% Processamento: mesas, telas e profissionais.
\foreach \x in {3.35,5.15}{%
\draw[ativo!70!black,line width=1pt] (\x,2.12)--++(1.15,0) (\x+.1,2.12)--++(0,-.32) (\x+1.05,2.12)--++(0,-.32);
\fill[ativo!60] (\x+.25,2.3) rectangle ++(.6,.37);\fill[white] (\x+.31,2.37) rectangle ++(.48,.23);
\draw[ativo] (\x+.55,2.3)--(\x+.55,2.12);
\fill[blue_theme] (\x+.1,2.65) circle(.075);\draw[blue_theme,thick] (\x+.1,2.56)--(\x+.1,2.27)--(\x+.32,2.21);}
\pessoaDados{7.05}{3.4}
% Térreo: energia e infraestrutura.
\fill[lacuna!70] (3.2,.5) rectangle (4.2,1.25);
\node[font=\large,blue_theme] at (3.7,.88){$+$};
\draw[blue_theme,line width=1pt] (4.45,.6)--(6.9,.6);
\foreach \x in {4.6,5.2,5.8,6.4}{\draw[blue_theme] (\x,.6)--(\x,1.05);}
\node[anchor=east,blue_theme,font=\scriptsize,align=right] at (2.4,3.75){2º andar\\\textbf{Memória}};
\node[anchor=east,ativo,font=\scriptsize,align=right] at (2.4,2.3){1º andar\\\textbf{Processamento}};
\node[anchor=east,black!65,font=\scriptsize,align=right] at (2.4,.85){Térreo\\Alimentação};
% Elevador de dados em uma coluna contínua.
\fill[white!60!blue_theme] (7.65,.35) rectangle (8.25,4.4);
\draw[blue_theme!40] (7.65,.35) rectangle (8.25,4.4);
\draw[<->,blue_theme,very thick] (7.95,1.95)--(7.95,3.9);
\fill[lacuna!75] (7.76,2.9) rectangle (8.14,3.24);\node[font=\tiny,blue_theme] at (7.95,3.07){01};
\node[font=\scriptsize,align=left,anchor=west] at (9.45,2.9){Conexões\\verticais:\\menos percurso,\\mais paralelismo};
\node[font=\scriptsize,blue_theme] at (5.6,5.18){Mesmo prédio: funções em andares especializados};
\fi
\path[use as bounding box] (-.3,-.3) rectangle (12,5.55);
\end{tikzpicture}}

\newcommand{\figura}[1]{\begin{center}\begin{adjustbox}{max width=\linewidth,max totalheight=5.1cm}#1\end{adjustbox}\end{center}}
\author[L. Nogueira / L. Soares]{\texorpdfstring{Luan F. Nogueira\\Luiz F. de Melo Soares}{Luan F. Nogueira e Luiz F. de Melo Soares}}
\institute[UFC]{\department{Minicurso de Microeletrônica}\\\ufc}
\date{SEELECT}

% Esboços vetoriais de inversores CMOS. Geometrias ilustrativas, sem DRC.
\newcommand{\inversorGeometria}[2]{%
% #1 = W em nm, #2 = L em nm. Apenas as dimensões do canal são proporcionais.
\begin{tikzpicture}[x=1cm,y=1cm,font=\small,line join=round]
\pgfmathsetmacro{\alturaW}{#1/150}
\pgfmathsetmacro{\larguraL}{#2/150}
\pgfmathsetmacro{\direita}{1.1+\larguraL}
\pgfmathsetmacro{\gatex}{.55+.5*\larguraL}
\pgfmathsetmacro{\contatod}{\direita-.22}
\fill[ativo!9,draw=ativo!35,dashed,rounded corners=3pt] (-.15,2.15) rectangle (\direita+.15,2.75+\alturaW);
\fill[blue_theme!9,draw=blue_theme!35,dashed,rounded corners=3pt] (-.15,-.15) rectangle (\direita+.15,.4+\alturaW);
\fill[ativo!40,draw=ativo!80!black] (0,2.45) rectangle (\direita,2.45+\alturaW);
\fill[blue_theme!40,draw=blue_theme] (0,.1) rectangle (\direita,.1+\alturaW);
% Gate comum cruza as duas regiões ativas, como no esboço de um inversor.
\fill[red!55,draw=red!75!black] (.55,-.25) rectangle (.55+\larguraL,2.75+\alturaW);
\foreach \y in {.1,2.45}{%
  \fill[black!45] (.15,\y+.5*\alturaW-.08) rectangle (.31,\y+.5*\alturaW+.08);
  \fill[black!45] (\contatod-.08,\y+.5*\alturaW-.08) rectangle (\contatod+.08,\y+.5*\alturaW+.08);}
% Metal de Source, Gate e Drain.
\draw[line width=1pt,black!65] (.23,2.45+.5*\alturaW)--(-.3,2.45+.5*\alturaW)--(-.3,4.75) node[above,font=\scriptsize]{$\VDD$};
\draw[line width=1pt,black!65] (.23,.1+.5*\alturaW)--(-.3,.1+.5*\alturaW)--(-.3,-.7) node[ground]{};
\draw[line width=1pt,black!65] (\contatod,2.45+.5*\alturaW)--(\direita+.35,2.45+.5*\alturaW)--(\direita+.35,.1+.5*\alturaW)--(\contatod,.1+.5*\alturaW);
\draw[line width=1pt,black!65] (\direita+.35,1.7)--(\direita+.7,1.7) node[right,font=\scriptsize]{$Y$};
\draw[line width=1pt,red!75!black] (\gatex,1.7)--(-.7,1.7) node[left,font=\scriptsize]{$A$};
\draw[<->,blue_theme] (-.78,.1)--node[left,font=\scriptsize]{$W$}(-.78,.1+\alturaW);
\draw[densely dotted,blue_theme] (.55,-.25)--(.55,-.55) (.55+\larguraL,-.25)--(.55+\larguraL,-.55);
\draw[<->,blue_theme] (.55,-.48)--node[below,font=\scriptsize]{$L$}(.55+\larguraL,-.48);
\path[use as bounding box] (-1.1,-1.05) rectangle (2.9,5.2);
\end{tikzpicture}}
\newcommand{\tresInversoresGeometria}{%
\begin{columns}[T,onlytextwidth]
\begin{column}{.32\textwidth}\centering
\textcolor{blue_theme}{\textbf{Referência}}\par
\begin{adjustbox}{max width=\linewidth,max totalheight=3.7cm}\inversorGeometria{120}{45}\end{adjustbox}\par
{\scriptsize $W=120$ nm; $L=45$ nm\par $W/L\approx2{,}67$}
\end{column}
\begin{column}{.32\textwidth}\centering
\textcolor{ativo}{\textbf{Dobrar W}}\par
\begin{adjustbox}{max width=\linewidth,max totalheight=3.7cm}\inversorGeometria{240}{45}\end{adjustbox}\par
{\scriptsize $W=240$ nm; $L=45$ nm\par $W/L\approx5{,}33$}
\end{column}
\begin{column}{.32\textwidth}\centering
\textcolor{blue_theme}{\textbf{Dobrar L}}\par
\begin{adjustbox}{max width=\linewidth,max totalheight=3.7cm}\inversorGeometria{120}{90}\end{adjustbox}\par
{\scriptsize $W=120$ nm; $L=90$ nm\par $W/L\approx1{,}33$}
\end{column}
\end{columns}}

\graphicspath{{../praticas/imagens/}{praticas/imagens/}}
\newcommand{\cod}[1]{\texttt{\detokenize{#1}}}
\newcommand{\acao}[2]{\textcolor{blue_theme}{\textbf{#1}}\par\smallskip #2\par\medskip}
\newcommand{\telaPasso}[4]{%
\begin{frame}{#1}
\begin{columns}[T,onlytextwidth]
\begin{column}{.59\textwidth}\centering\includegraphics[width=\linewidth,height=4.5cm,keepaspectratio]{#2}\end{column}
\begin{column}{.38\textwidth}{\small #3}\end{column}
\end{columns}
\fonte{\textbf{Conferir:} #4}
\end{frame}}
\newcommand{\telaAmpla}[4]{%
\begin{frame}{#1}
\begin{center}\includegraphics[width=.97\linewidth,height=4.1cm,keepaspectratio]{#2}\end{center}
{\small #3\par}
\fonte{#4}
\end{frame}}
\newcommand{\duasTelas}[6]{%
\begin{frame}{#1}
\begin{columns}[T,onlytextwidth]
\begin{column}{.48\textwidth}\centering\includegraphics[width=\linewidth,height=3.8cm,keepaspectratio]{#2}\par\smallskip{\scriptsize #4}\end{column}
\begin{column}{.48\textwidth}\centering\includegraphics[width=\linewidth,height=3.8cm,keepaspectratio]{#3}\par\smallskip{\scriptsize #5}\end{column}
\end{columns}
\fonte{\textbf{Conferir:} #6}
\end{frame}}
\title[Práticas • Virtuoso]{\huge\textbf{Do transistor ao chip}}
\begin{document}
\capa{Roteiro visual de práticas · Cadence Virtuoso}{NMOS/PMOS • portas CMOS • layout e DRC/LVS}
\section{Percurso}
\begin{frame}{Três módulos para acompanhar na bancada}
\centering\small\begin{tabular}{lll}\toprule
Módulo & O que construir & Tempo reservado\\\midrule
\hyperlink{pratica1}{\textbf{Aula 1}} & Testes DC de NMOS e PMOS & 50 min\\[2mm]
\hyperlink{pratica2}{\textbf{Aula 2}} & Inversor, NAND e NOR & 60 min\\[2mm]
\hyperlink{pratica3}{\textbf{Aula 3}} & Layout, DRC e LVS do inversor & 95 min\\\bottomrule
\end{tabular}\par\bigskip
\textbf{Ambiente das capturas}\par
Virtuoso Studio IC23.1 · \cod{gpdk045} · Spectre\par
ADE Assembler · Layout XL · Assura
\fonte{Adaptação de praticas/QUICK\_START.md, incorporado ao projeto. Os nomes de menu permanecem como na interface. Source, Gate e Drain nomeiam os terminais MOS. A pausa de 5 minutos está nas apresentações, antes de cada prática.}
\end{frame}
% Introdução curta, compartilhada entre a aula 1 e o roteiro de bancada.
\begingroup
\begin{frame}{Nossa tecnologia na prática: gpdk045}
\begin{center}\begin{tikzpicture}[x=1cm,y=1cm,font=\small]
\draw[blue_theme,thick,fill=blue_theme!6] (3.7,0) rectangle (6.4,2.2);
\node[blue_theme,font=\bfseries] at (5.05,1.35){gpdk045};\node[font=\scriptsize] at (5.05,.65){CMOS genérico · 45 nm};
\node[align=center,font=\scriptsize] (a) at (1.1,3){\textbf{Schematic}\\Símbolos e terminais};
\node[align=center,font=\scriptsize] (b) at (9,3){\textbf{Simulação}\\Modelos elétricos};
\node[align=center,font=\scriptsize] (c) at (1.1,-.55){\textbf{Layout}\\Células parametrizadas};
\node[align=center,font=\scriptsize] (d) at (9,-.55){\textbf{Verificação}\\Regras DRC/LVS};
\draw[seta] (3.7,1.7)--(a);\draw[seta] (6.4,1.7)--(b);\draw[seta] (3.7,.3)--(c);\draw[seta] (6.4,.3)--(d);
\path[use as bounding box] (-.5,-1.15) rectangle (10.6,3.65);
\end{tikzpicture}\end{center}
{\small\centering Na nossa montagem: $V_{DD}=1$ V, $W=120$ nm, $L=45$ nm.\par
\texttt{nmos1v} / \texttt{pmos1v} · um finger · multiplicador 1.\par}
\fonte{PDK = Process Design Kit. Fonte: \href{https://community.cadence.com/cadence_technology_forums/f/custom-ic-design/36796/basics-of-pdk/1349734}{Cadence — Basics of PDK}: kit genérico para ensino e demonstração. Parâmetros aplicados: Quick Start da prática; 1 V é a alimentação usada neste minicurso.}
\end{frame}

\begin{frame}{O mesmo inversor CMOS, três geometrias}
\tresInversoresGeometria
\fonte{Vista superior conceitual: PMOS em verde, NMOS em azul, Gate em vermelho e contatos/metal em cinza. As setas W/L indicam o canal. Os valores valem para ambos os transistores em cada exemplo, com um finger e multiplicador 1. Contatos, poços e espaçamentos simplificados; não é um layout fabricável. As duas variações são propostas de comparação, sem novas curvas simuladas.}
\end{frame}

\begin{frame}{Como W/L aparece na corrente de Drain?}
\small\centering
\textcolor{blue_theme}{\textbf{NMOS em saturação · aproximação de canal longo}}\par\medskip
\[
 I_{DS}\approx\frac{1}{2}\,\mu_n C_{\mathrm{ox}}\,
 \textcolor{blue_theme}{\frac{W}{L}}\,(V_{GS}-V_{TN})^2
\]
{\scriptsize $V_{GS}>V_{TN}$ e $V_{DS}\geq V_{GS}-V_{TN}$; corpo ligado ao Source.}\par\medskip
$\mu_n$: mobilidade\qquad $C_{\mathrm{ox}}$: capacitância por área\par
$V_{TN}$: limiar\qquad $W/L$: razão geométrica\par\medskip
\begin{tabular}{lccc}\toprule
Modelo ideal & Referência & Dobrar W & Dobrar L\\\midrule
Corrente relativa & $1\times$ & $2\times$ & $0{,}5\times$\\\bottomrule
\end{tabular}
\fonte{Mesmas tensões, mobilidade, Cox e limiar. Para o módulo da corrente PMOS, usar $\mu_p$, $V_{SG}$ e $|V_{TP}|$. Fonte: \href{https://web.mit.edu/6.101/www/s2020/handouts/L06.pdf}{MIT 6.101, aula 6, pp. 19–20}. A lei quadrática explica tendências; a simulação em 45 nm usa os modelos do PDK e não precisa seguir exatamente essas proporções. W maior também aumenta capacitâncias.}
\end{frame}
\endgroup

\section{Prática 1}
\telaPasso{A1 · 01 — Criar a biblioteca}{01-nova-biblioteca.png}{%
\hypertarget{pratica1}{}%
\acao{Library Manager}{File → New → Library}
\acao{Nome}{\cod{minicurso_cmos}}
\acao{Tecnologia}{Vincular à tecnologia existente \cod{gpdk045}.}
}{A captura mostra o menu. A janela de vinculação tecnológica não foi registrada.}

\duasTelas{A1 · 02 — Criar a célula de teste}{02-nova-celula.png}{03-configuracao-celula.png}{%
File → New → Cell View}{%
Library: \cod{minicurso_cmos}\par Cell: \cod{a1_tb_nmos}\par View/Type: \cod{schematic}\par Open with: Schematics L
}{Confirmar os campos e clicar em OK.}

\telaAmpla{A1 · 03 — Reconhecer o editor}{04-schematic-vazio.png}{%
Confira no título: biblioteca \cod{minicurso_cmos}, célula \cod{a1_tb_nmos} e vista \cod{schematic}.
}{Com o foco na área de desenho, a tecla i abre Add Instance neste ambiente; Esc encerra a inserção.}

\telaPasso{A1 · 04 — Inserir o NMOS}{05-nmos1v.png}{%
\acao{Add Instance}{\cod{gpdk045/nmos1v}\par View: \cod{symbol}}
\acao{Modelo}{\cod{g45n1svt}}
\acao{Dimensões iniciais}{$L=45$ nm; $W=120$ nm.\par Um finger; multiplicador 1.}
}{Gate à esquerda, Drain acima, Source abaixo e corpo à direita, na orientação registrada.}

\duasTelas{A1 · 05 — Fontes e referências}{06-fonte-vdc.png}{07-simbolo-vdd.png}{%
\cod{analogLib/vdc}: duas fontes independentes para Gate e Drain.}{%
\cod{analogLib/vdd}: identifica a rede; a tensão vem de uma fonte. Inserir também \cod{analogLib/gnd}.
}{Os símbolos vdd/gnd nomeiam redes; não aplicam tensão por si sós. A seleção específica de gnd não foi capturada.}

\telaAmpla{A1 · 06 — Conectar o teste do NMOS}{11-montagem-nmos-com-referencia-vdd.png}{%
Source e corpo → GND.\quad Gate → positivo de V0.\quad Drain → positivo de V3.
}{V0 = vgs; V3 = vds. Negativos das fontes em GND. V2 = 1 V explica a rede vdd, mas não está ligada a Gate ou Drain neste teste.}

\telaPasso{A1 · 07 — Conferir a fonte de alimentação}{09-propriedades-vdc.png}{%
\acao{Instância mostrada}{\cod{V2}}
\acao{DC voltage}{\cod{1} = 1 V}
\acao{Demais fontes}{Em V0 usar \cod{vgs}; em V3 usar \cod{vds}.}
}{A janela é de V2, não das fontes parametrizadas V0 e V3. Conferir também seus negativos em GND.}

\telaPasso{A1 · 08 — Conferir o dispositivo}{10-propriedades-nmos.png}{%
\acao{Instância}{\cod{NM0}}
{\scriptsize\begin{tabular}{ll}Modelo & \cod{g45n1svt}\\Length & \cod{45n}\\Total Width & \cod{120n}\\Fingers & 1\\Multiplier & 1\end{tabular}}\par\medskip
\textbf{Check and Save}\par Examinar conexões e avisos.
}{Bodytie Type = None não substitui a conexão explícita do corpo ao Source/GND. Proximidade gráfica não garante ligação elétrica.}

\duasTelas{A1 · 09 — Abrir o ADE Assembler}{12-launch-ade.png}{13-inicio-assembler.png}{%
No schematic: Launch → ADE Assembler.}{%
A sessão abre com a aba maestro e o painel Data View.
}{O fluxo registrado é do Assembler. O Explorer aparece no menu, mas não tem um passo a passo separado neste material.}

\duasTelas{A1 · 10 — Criar e conferir o teste}{14-data-view.png}{15-primeiro-teste.png}{%
Na seção Tests, definir um teste para \cod{minicurso_cmos/a1_tb_nmos}.}{%
Simulator: \cod{spectre}.\par Conferir Analyses e Design Variables no teste criado.
}{A janela intermediária de seleção do design não foi capturada. Use variáveis do próprio teste nesta caracterização.}

\telaAmpla{A1 · 11 — Parametrizar VGS}{19-variaveis-varredura-vgs.png}{%
\cod{vgs}: From/To; Step Type = Auto; From = \cod{100m}; To = \cod{700m}; Total Steps = 5.
}{Valores registrados na família: 0,10; 0,25; 0,40; 0,55; 0,70 V. A referência vds = 500m é substituída pela varredura DC.}

\telaPasso{A1 · 12 — Varrer VDS em DC}{16-analise-dc-vds.png}{%
{\scriptsize\begin{tabular}{@{}ll@{}}Análise & \cod{dc}\\Sweep Variable & Design Variable\\Variable Name & \cod{vds}\\Range & Start-Stop\\Start / Stop & 0 / 1 V\\Sweep Type & Automatic\end{tabular}}\par\medskip
Marcar Save DC Operating Point e Enabled; confirmar em OK.
}{Para cada VGS, a análise percorre VDS de 0 a 1 V. O título deste formulário mostra ADE Explorer, dentro do fluxo registrado no Assembler.}

\telaPasso{A1 · 13 — Selecionar corrente e executar}{17-saidas-e-execucao.png}{%
\acao{Saída}{Corrente no terminal Drain de NM0:\par\cod{/NM0/D}}
\acao{Conferir antes de executar}{Análise DC, variáveis e configuração dos modelos do PDK.}
\acao{Executar}{Botão triangular verde. Salvar a sessão.}
}{Selecionar um terminal mede corrente; selecionar uma rede mede tensão. A configuração dos modelos e a seleção no schematic ainda não têm captura.}

\telaAmpla{A1 · 14 — Ler as curvas do NMOS}{18-curvas-id-vds.png}{%
Horizontal: VDS, 0–1 V.\quad Vertical: corrente em Drain, em $\mu$A.\par
Cada curva mantém VGS fixo; elevar VGS aumenta a condução.
}{A curva de VGS = 0,70 V chega a aproximadamente 10 µA em VDS = 1 V. Não exigir um patamar horizontal nem extrair o limiar apenas desta figura.}

\telaAmpla{A1 · 15 — Montar o teste do PMOS}{20-montagem-pmos.png}{%
\cod{a1_tb_pmos}: Source e corpo em VDD = 1 V.\par
$V_{SG}=V_S-V_G$ e $V_{SD}=V_S-V_D$: positivos das fontes no Source.
}{PM0: g45p1svt, L = 45 nm, W = 120 nm, multiplicador 1. Negativos das fontes de controle em Gate e Drain, respectivamente.}

\telaAmpla{A1 · 16 — Comparar as curvas do PMOS}{21-curvas-pmos.png}{%
Varrer \cod{vsd} de 0 a 1 V e repetir para \cod{vsg}: 0,10; 0,25; 0,40; 0,55; 0,70 V.
}{A saída mostrada é /PM0/S, corrente em Source, e não /PM0/D. Comparar polaridades e terminal medido; W e L iguais não garantem correntes iguais.}

\section{Prática 2}
\telaPasso{A2 · 01 — Montar o inversor}{22-montagem-inversor.png}{%
\hypertarget{pratica2}{}%
\acao{Célula}{\cod{a2_inv}}
\acao{Entrada vin}{Gates unidos.}
\acao{Saída vout}{Drains unidos.}
\acao{Alimentação}{PMOS: Source/corpo em VDD. NMOS: Source/corpo em GND.}
}{W = 120 nm; L = 45 nm; multiplicador 1; VDD = 1 V. A fonte mostrada está no schematic: evitar duplicá-la no testbench.}

\duasTelas{A2 · 02 — Criar a vista symbol}{23-menu-criar-symbol.png}{24-cellview-symbol.png}{%
Create → Cellview → From Cellview…}{%
\cod{a2_inv}: From View = schematic; To View = symbol; Tool = schematicSymbol.
}{Conferir biblioteca e célula; manter Display Cellview e Edit Options marcados.}

\duasTelas{A2 · 03 — Organizar os pinos e o desenho}{25-opcoes-symbol.png}{26-symbol-inversor.png}{%
Left Pins: \cod{vin}.\par Right Pins: \cod{vout}.\par Top/Bottom Pins vazios nesta versão.}{%
Ajustar o desenho, preservando terminais, pontos de conexão e rótulos da célula/instância.
}{O símbolo representa a função; o comportamento elétrico continua no schematic. A personalização gráfica não foi capturada passo a passo.}

\duasTelas{A2 · 04 — Conferir as duas vistas}{27-vista-symbol.png}{28-vista-schematic.png}{%
A vista \cod{symbol} será instanciada no teste.}{%
A vista \cod{schematic} mantém os transistores e suas conexões.
}{Ambas pertencem a \cod{minicurso_cmos/a2_inv}. Nesta versão, o símbolo expõe apenas vin/vout; a alimentação interna continua necessária.}

\telaPasso{A2 · 05 — Montar o testbench do inversor}{29-testbench-inversor.png}{%
\acao{Célula}{\cod{a2_tb_inv}}
\acao{Instância I0}{\cod{a2_inv/symbol}}
\acao{Entrada}{V0 entre vin e GND.}
\acao{Saída}{Selecionar a tensão de \cod{/vout} no ADE.}
}{Conferir a alimentação interna. Os parâmetros completos de pulso do inversor não foram registrados; o desenho da fonte não confirma seus tempos.}

\telaAmpla{A2 · 06 — Conferir a inversão no tempo}{30-transiente-inversor.png}{%
\textcolor{purple}{Roxo: /vout.}\quad\textcolor{orange}{Laranja: /vin.}\par
Nos intervalos estáveis: entrada alta → saída baixa; entrada baixa → saída alta.
}{Resultado registrado entre 0 e 40 ns. O objetivo é a função lógica; curva de transferência DC do inversor e investigação dos picos ficam fora desta aula.}

\telaPasso{A2 · 07 — Construir a NAND}{31-schematic-nand.png}{%
\acao{Célula}{\cod{a2_nand2}}
\acao{Subida}{PMOS em paralelo.}
\acao{Descida}{NMOS em série.}
\acao{A, B e Y}{Manter exatamente os mesmos nomes no circuito, símbolo e teste.}
}{Y é o nó comum dos Drains PMOS e do Drain do NMOS superior. O nó entre NMOS é interno. Corpos NMOS em GND e PMOS em VDD, inclusive nas pilhas.}

\duasTelas{A2 · 08 — Símbolo e teste da NAND}{32-symbol-nand.png}{37-testbench-nand.png}{%
Entradas A/B à esquerda; saída Y à direita. Preservar os pontos de conexão.}{%
\cod{a2_tb_nand2}: duas fontes vpulse, com negativos em GND; medir A, B e Y.
}{Conferir alimentação de 1 V e executar Check and Save. W = 120 nm, L = 45 nm, multiplicador 1 por transistor.}

\telaPasso{A2 · 09 — Construir a NOR}{34-schematic-nor.png}{%
\acao{Célula}{\cod{a2_nor2}}
\acao{Subida}{PMOS em série.}
\acao{Descida}{NMOS em paralelo.}
\acao{Corpos}{Todos os PMOS em VDD; todos os NMOS em GND.}
}{Y une o Drain do PMOS inferior aos Drains NMOS. O nó entre PMOS é interno. Os símbolos de alimentação não substituem a fonte de 1 V.}

\duasTelas{A2 · 10 — Símbolo e teste da NOR}{35-symbol-nor.png}{38-testbench-nor.png}{%
Repetir a geração de symbol com A/B à esquerda e Y à direita.}{%
\cod{a2_tb_nor2}: aplicar os mesmos quatro estados usados na NAND e medir A, B e Y.
}{Conferir os nomes das redes, a alimentação e as ligações de corpo. Executar Check and Save antes do ADE.}

\begin{frame}{A2 · 11 — Configurar os estímulos e tran}
\small\centering
\begin{tabular}{lll}\toprule
Campo & Entrada A & Entrada B\\\midrule
Voltage 1 / Voltage 2 & 0 / 1 V & 0 / 1 V\\
Delay time & 0 & 0\\
Pulse Width & 5 ns & 10 ns\\
Period & 10 ns & 20 ns\\
Rise time na NAND & 100 ps & 100 ps\\
Rise time na NOR & 20 ps & 20 ps\\\bottomrule
\end{tabular}\par\bigskip
\textbf{Análise:} tran\qquad\textbf{Stop Time:} \cod{40n}\par\medskip
A alterna mais rápido que B; observar as quatro combinações.
\fonte{Delay = 0 e Stop Time = 40n foram informados pelo responsável; tempos de subida aparecem nos testbenches. Para descida, 100 ps é a referência sugerida no roteiro, sem captura de confirmação. A configuração completa do pulso do inversor permanece não registrada.}
\end{frame}

\telaPasso{A2 · 12 — Validar a NAND nos intervalos estáveis}{33-transiente-nand.png}{%
{\scriptsize\centering\begin{tabular}{rccc}\toprule $t$ & A&B&Y\\\midrule
2 ns&1&1&0\\7 ns&0&1&1\\12 ns&1&0&1\\17 ns&0&0&1\\\bottomrule\end{tabular}\par}\medskip
\textcolor{red}{A: vermelho.}\par\textcolor{ativo}{B: verde.}\par\textcolor{purple}{Y: magenta.}\par\medskip
$Y=\overline{AB}$
}{Ler longe das bordas, após o assentamento. Glitches e picos não são objeto de investigação nesta atividade.}

\telaPasso{A2 · 13 — Validar a NOR nos intervalos estáveis}{36-transiente-nor.png}{%
{\scriptsize\centering\begin{tabular}{rccc}\toprule $t$ & A&B&Y\\\midrule
2 ns&1&1&0\\7 ns&0&1&0\\12 ns&1&0&0\\17 ns&0&0&1\\\bottomrule\end{tabular}\par}\medskip
\textcolor{cyan!70!black}{A: ciano.}\par\textcolor{purple}{B: roxo.}\par\textcolor{orange}{Y: laranja.}\par\medskip
$Y=\overline{A+B}$
}{Saída alta somente em 00. Os gráficos registrados verificam a lógica das três portas; a próxima aula leva o inversor ao layout.}

\section{Prática 3}
\telaPasso{A3 · 01 — Preparar o inversor e abrir Layout XL}{39-launch-layout-xl.png}{%
\hypertarget{pratica3}{}%
\acao{Célula}{\cod{a3_inv}: um PMOS e um NMOS, W = 120 nm, L = 45 nm.}
\acao{Alimentação}{Fonte ideal no testbench. Usar as redes herdadas \cod{vdd!} e \cod{gnd!} da montagem.}
\acao{Abrir}{Launch → Layout XL}
}{O schematic usa \cod{vdd_inherit} e \cod{gnd_inherit}. Preservar esses nomes ao preparar terminais físicos e LVS, sem trocá-los silenciosamente por redes locais.}

\duasTelas{A3 · 02 — Criar a vista layout}{40-layout-startup.png}{41-nova-vista-layout.png}{%
Startup Option:\par Layout = Create New;\par Configuration = Automatic.}{%
Library: \cod{minicurso_cmos}\par Cell: \cod{a3_inv}\par View/Type: \cod{layout}\par Open with: Layout XL
}{O layout é uma nova vista da mesma célula. Confirmar os campos e clicar em OK.}

\telaAmpla{A3 · 03 — Reconhecer a área de layout}{42-layout-vazio-schematic.png}{%
À esquerda, camadas do PDK; à direita, schematic de referência.\par
A vista ainda está vazia: criar a vista não gera um layout concluído.
}{Próxima operação: gerar os dispositivos a partir do schematic no Layout XL, posicionar PMOS acima e NMOS abaixo e iniciar as conexões. As janelas intermediárias não foram registradas.}

\telaPasso{A3 · 04 — Relacionar o layout atualizado}{43-layout-inversor.png}{%
\acao{PMOS superior}{Source e corpo em \cod{vdd!}.}
\acao{NMOS inferior}{Source e corpo em \cod{gnd!}.}
\acao{vin}{Gates dos dois dispositivos.}
\acao{vout}{Drains dos dois dispositivos.}
}{Comparar com \cod{a3_inv/schematic} e conferir camadas, contatos e nomes das redes. Esta é a captura atualizada da montagem; os resultados das verificações aparecem nas próximas etapas.}

\begin{frame}{A3 · 05 — Configurar a tecnologia no Assura}
\small
\textcolor{blue_theme}{\textbf{Se a tecnologia aparecer como -undefined-}}\par\medskip
Na CIW, consultar o caminho da biblioteca:\par\smallskip
\cod{ddGetObj("gpdk045")~>readPath}\par\medskip
Em \textbf{Assura → Technology…}, selecionar:\par\smallskip
{\scriptsize\cod{/home/tools/cadence/gpdk/gpdk045_v_6_0/}\par\cod{assura_tech.lib}}\par\medskip
\textbf{Tecnologia esperada:} \cod{gpdk045_av}\par
\textbf{Rule Set:} \cod{default}
\fonte{O caminho é específico do servidor do minicurso; em outra instalação, localizar o arquivo correspondente no PDK. Reconhecer as regras é uma configuração anterior à verificação do layout.}
\end{frame}

\duasTelas{A3 · 06 — Localizar o DRC no Assura}{47-menu-assura-drc.png}{44-configuracao-drc-assura.png}{%
No layout: Assura → Run DRC…}{%
Configuração de referência:\par gpdk045\_av · default\par Regras: assuraDRC.rul
}{As regras precisam estar reconhecidas antes da execução. O próximo formulário registra a versão final usada na verificação.}

\telaPasso{A3 · 07 — Executar DRC na versão final}{48-formulario-drc-final.png}{%
\acao{Design}{\cod{a3_inv/layout}\par Área: Full}
\acao{Regras}{\cod{gpdk045_av}\par default · \cod{assuraDRC.rul}}
\acao{Executar}{Conferir nome e diretório na sessão; preservar os resultados.}
}{A captura tem Run Name vazio e Run Directory = ponto. Escolher dados de execução válidos no ambiente; o formulário sozinho não é o resultado do DRC.}

\telaPasso{A3 · 08 — Resultado registrado: DRC sem erros}{46-drc-sem-erros.png}{%
\acao{Mensagem do Assura}{\textbf{No DRC errors found.}}
\acao{O que ela confirma}{As regras de DRC executadas não encontraram violações nesta versão do layout.}
\acao{Próximo passo}{Conferir a correspondência elétrica pelo LVS.}
}{Se o layout mudar durante uma correção, repetir o DRC na versão alterada. A captura é o resultado fornecido para o minicurso.}

\duasTelas{A3 · 09 — Localizar o LVS no Assura}{49-menu-assura-lvs.png}{45-configuracao-lvs-assura.png}{%
No layout: Assura → Run LVS…}{%
Configuração de referência:\par extract.rul · compare.rul\par LVSinclude.rsf
}{Comparar schematic e layout da mesma célula \cod{a3_inv}. Tecnologia gpdk045\_av, Rule Set default.}

\telaPasso{A3 · 10 — Executar a comparação final de LVS}{50-formulario-lvs-final.png}{%
\acao{Vistas de a3\_inv}{\cod{schematic}\par\cod{layout}}
\acao{Regras}{\cod{extract.rul}\par\cod{compare.rul}\par\cod{LVSinclude.rsf}}
\acao{Executar}{Conferir nome/diretório e ler o resumo.}
}{Biblioteca \cod{minicurso_cmos}, tecnologia gpdk045\_av e Rule Set default. O resultado mostrado a seguir identifica a execução como \cod{a3_inv}.}

\telaPasso{A3 · 11 — Resultado registrado: LVS corresponde}{51-lvs-match.png}{%
\acao{Mensagem do Assura}{\textbf{Schematic and Layout Match.}}
\acao{Resumo}{Zero divergências de redes, dispositivos, pinos e parâmetros.}
\acao{Navegação}{Yes permite visualizar os resultados da nova execução.}
}{A janela também indica Total DRC violations: 0. A pergunta de navegação não é uma falha de LVS; a correspondência já está informada no resumo.}

\begin{frame}{A3 · 12 — Concluir com as evidências de DRC e LVS}
\figura{\begin{tikzpicture}[node distance=12mm,font=\small]
\node[draw=blue_theme,align=center,minimum width=2.5cm,minimum height=1.2cm] (a){\textbf{DRC}\\Sem erros};
\node[draw=blue_theme,right=of a,align=center,minimum width=2.5cm,minimum height=1.2cm] (b){\textbf{LVS}\\Schematic/Layout Match};
\node[draw=blue_theme,right=of b,align=center,minimum width=2.5cm,minimum height=1.2cm] (c){\textbf{Registrar}\\Preservar resultados};
\draw[seta] (a)--(b);\draw[seta] (b)--(c);
\node[below=10mm of b,font=\scriptsize,align=center] (d){Alterou o layout?\\Repetir as verificações pertinentes.};
\draw[seta] (d.west)-|(a.south);
\end{tikzpicture}}
\chamada{O layout verificado corresponde ao schematic do inversor.}
\fonte{Conclusão apoiada nas capturas de DRC e LVS fornecidas. A aula termina nessas verificações; extração de parasitas e simulação após extração ficam fora do escopo. Retomar o encerramento e o roadmap.}
\end{frame}

\section{Créditos}
\begin{frame}{Referências e créditos}
\small\begin{itemize}
\item Desenvolvimento: \textbf{Luan F. Nogueira e Luiz F. de Melo Soares}.
\item Fonte das etapas, parâmetros e capturas reais: \textit{Quick Start — Práticas de CMOS no Cadence Virtuoso}, fornecido para este minicurso, edição incorporada em 31/10/2026.
\item As imagens são as capturas originais de \cod{praticas/imagens/}; não são telas geradas ou resultados produzidos por esta adaptação.
\item PDK e dimensionamento: Cadence, \href{https://web.mit.edu/6.101/www/s2020/handouts/L06.pdf}{MIT 6.101 — MOSFETs}; fontes e hipóteses nos três quadros introdutórios.
\item Template Beamer UFC: Maurício Moreira Neto. Modelo original 4:3.
\end{itemize}
\fonte{O roteiro textual completo acompanha o pacote de entrega. Parâmetros informados sem captura são identificados nas etapas; os resultados de DRC/LVS reproduzem as capturas fornecidas.}
\end{frame}
\end{document}
